%==============================================================
%  lecons/analyse/rolle.tex — Théorème de Rolle
%==============================================================

\lecontitle{Théorème de Rolle}{Analyse réelle — Fonctions dérivables}

%=============================================================
%  § 1 — ÉNONCÉ
%=============================================================
\secnum{navyblue}{1}{Énoncé et signification}

\begin{tcolorbox}[thmbox]
  \textbf{Théorème (Rolle, 1691).}\enspace%
  \index{Rolle (théorème de)}
  \textit{Soit $f : [a,b] \to \mathbb{R}$ une fonction vérifiant :}
  \begin{enumerate}
    \item \textit{$f$ est \textbf{continue} sur $[a,b]$ ;}
    \item \textit{$f$ est \textbf{dérivable} sur $]a,b[$ ;}
    \item \textit{$f(a) = f(b)$.}
  \end{enumerate}
  \textit{Alors il existe $c \in\, ]a,b[$ tel que $f'(c) = 0$.}
\end{tcolorbox}

\rubriq{Signification géométrique}
Si la courbe $y = f(x)$ repart du même niveau qu'elle a quitté, elle possède
au moins une tangente horizontale quelque part entre $a$ et $b$.
C'est une propriété de \emph{valeur intermédiaire pour la dérivée}.

\rubriq{Nécessité des trois hypothèses}
Les trois conditions sont indépendantes et toutes nécessaires (voir §3).

\decsep{}

%=============================================================
%  § 2 — DÉMONSTRATION
%=============================================================
\secnum{navyblue}{2}{Démonstration}

\begin{ideepreuve}
$f$ est continue sur un compact $[a,b]$, donc elle y atteint son maximum et son
minimum. Si l'un d'eux est atteint en un point intérieur $c$, alors $f'(c) = 0$.
L'hypothèse $f(a)=f(b)$ force au moins un extremum à être intérieur dès que $f$
n'est pas constante.
\end{ideepreuve}

\begin{mdframed}[style=proofbar]
  \textit{Cas 1 : $f$ est constante.}\enspace%
  Alors $f' \equiv 0$ sur $]a,b[$, et tout point de $]a,b[$ convient.

  \medskip
  \textit{Cas 2 : $f$ n'est pas constante.}\enspace%
  $f$ est continue sur le compact $[a,b]$, donc par le théorème des bornes
  atteintes\index{bornes atteintes (théorème des)}, elle y atteint son maximum $M$
  et son minimum $m$. Puisque $f$ n'est pas constante, $m < M$.
  Or $f(a) = f(b)$, donc $M$ et $m$ ne peuvent pas être atteints \emph{tous les deux}
  uniquement aux extrémités $a$ et $b$. Ainsi, au moins l'un d'eux est atteint en
  un point $c \in\, ]a,b[$.

  En ce point, $f$ admet un extremum local, donc par le théorème de Fermat
  \index{Fermat (théorème de)}, $f'(c) = 0$. \hfill$\blacksquare$
\end{mdframed}

\rubriq{Rappel — Théorème de Fermat}

\begin{mdframed}[style=proofbar]
  Si $f$ est dérivable en $c$ et admet un extremum local en $c$, alors $f'(c)=0$.
  \textit{Preuve :} Si $c$ est un maximum local, alors pour $h>0$ assez petit,
  $\frac{f(c+h)-f(c)}{h}\leq 0$ et $\frac{f(c-h)-f(c)}{-h}\geq 0$.
  En passant à la limite, $f'(c)\leq 0$ et $f'(c)\geq 0$, donc $f'(c)=0$. \hfill$\blacksquare$
\end{mdframed}

\decsep{}

%=============================================================
%  § 3 — APPLICATIONS, CONTRE-EXEMPLES, EXERCICES
%=============================================================
\secnum{exgreen}{3}{Applications, contre-exemples et exercices}

\rubriq{Applications}

\begin{mdframed}[style=exbar]
  \textbf{1. Unicité des racines.}\enspace%
  Si $f'$ ne s'annule pas sur $]a,b[$, alors $f$ admet \emph{au plus} une racine dans
  $]a,b[$. (Par l'absurde : deux racines $\alpha < \beta$ donneraient via Rolle
  un $c \in\,]\alpha,\beta[$ avec $f'(c)=0$.)

  \smallskip\noindent
  \textbf{2. Séparation des racines.}\enspace%
  Entre deux racines consécutives de $f'$, la dérivée $f'$ garde un signe
  constant (théorème de Darboux), donc $f$ y est monotone.
  Entre deux racines consécutives de $f$, $f'$ s'annule au moins une fois.

  \smallskip\noindent
  \textbf{3. Version généralisée (Rolle itéré).}\enspace%
  Si $f$ est $n$ fois dérivable et admet $n+1$ racines distinctes sur $[a,b]$,
  alors $f^{(n)}$ admet au moins une racine sur $]a,b[$.
  (Par récurrence, en appliquant Rolle $n$ fois.)
\end{mdframed}

\vspace{6pt}
\rubriq{Contre-exemples — nécessité des hypothèses}

\begin{mdframed}[style=cexbar]
  \textbf{Continuité indispensable.}\enspace%
  $f(x) = \begin{cases}x & x\in[0,1[\\ 0 & x=1\end{cases}$
  vérifie $f(0)=f(1)=0$ et $f'(x)=1\neq 0$ sur $]0,1[$. La discontinuité en $1$
  invalide le théorème.

  \smallskip\noindent
  \textbf{Dérivabilité indispensable.}\enspace%
  $f(x) = |x|$ sur $[-1,1]$ : $f(-1)=f(1)=1$, $f$ continue, mais
  $f'(x) = \operatorname{sgn}(x) \neq 0$ sur $]-1,1[\,\setminus\{0\}$
  et $f$ n'est pas dérivable en $0$.

  \smallskip\noindent
  \textbf{$f(a)=f(b)$ indispensable.}\enspace%
  $f(x) = x$ sur $[0,1]$ : $f'(x)=1\neq 0$ partout, car $f(0)\neq f(1)$.
\end{mdframed}

\vspace{6pt}
\exolabel

\begin{mdframed}[style=exobar]
  \begin{enumerate}
    \item Montrer que l'équation $x^5 - 5x + 1 = 0$ admet exactement trois
          racines réelles. (Étudier $f' = 5x^4 - 5 = 5(x^2-1)(x^2+1)$.)

    \item Soit $f$ dérivable sur $\mathbb{R}$ avec $f(0)=0$ et $f(1)=1$.
          Montrer qu'il existe $a,b \in\,]0,1[$, $a\neq b$, tels que
          $f'(a)+f'(b)=2$. (\textit{Indication :} appliquer le théorème des
          accroissements finis (corollaire de Rolle) à $f$ sur $[0,1/2]$ et
          $[1/2,1]$ ; la somme des deux pentes vaut $2$ indépendamment de $f(1/2)$.)

    \item (Rolle itéré) Soit $P$ un polynôme de degré $n$ ayant $n$ racines
          réelles distinctes. Montrer que $P'$ a exactement $n-1$ racines réelles,
          $P''$ exactement $n-2$, etc.

    \item Soit $f$ continue sur $[0,1]$, $\mathcal{C}^1$ sur $]0,1[$, avec
          $f(0)=f(1)=0$ et $f'$ monotone. Montrer que $f$ garde un signe constant
          sur $[0,1]$.
  \end{enumerate}
\end{mdframed}

\leconfooter{M.\ Rolle, \textit{Méthode pour résoudre les égalitez} (1691)}
